% discussion.tex -- Discussion
\section{Discussion}

Changes in investor decision rules spread selling through the market,
while trader resources and the dealer's reserve curve determined the
associated price decline. A majority entered Sell after enough traders
received a new response to their neighbors, even though other traders
retained their original rules (Fig.~\ref{fig:rule-quench}). The changed
response remained available whenever the relevant neighbor configuration
recurred, rather than requiring selected traders to remain in Sell.
Comparing a persistent rule change with a one-time instruction to sell
separated this continuing influence from the immediate change in current
decisions (Supplementary Table~\ref{tab:revision-pulse}). The alternative
rule screen further tested which responses could produce a majority,
so the selected rule was examined within a defined set of local changes
rather than treated as a universal description of investor behavior
(Supplementary Table~\ref{tab:revision-rules}). The paired design made
the no-switch path the comparison for the intervention at each starting
market~\cite{glasserman1992guidelines,law2015simulation}. Repeating the
pair across starting markets then measured how frequently a majority
formed within the observation window. This interpretation concerned
assigned decision rules and their propagation through neighbors; traders
did not learn the new rule, choose it for its profitability, or pass
the rule itself to other traders.

The decision instability could develop without a price response entering
later decisions. Holding the technical traders removed that route of
influence, yet the rule change still produced Sell majorities
(Fig.~\ref{fig:reservoir}). In the spatial examples, growing groups of
sellers linked local responses to changes across the grid
(Fig.~\ref{fig:dose-panels}). Such patterns resemble the growth of
ordered regions in other interacting systems away from
equilibrium~\cite{hinrichsen2000nonequilibrium,bray1994theory}, while
dependence on the changed fraction connects the model to accounts of
how local changes spread through a
population~\cite{granovetter1978threshold,watts2002simple,aizenman1988metastability,holroyd2003sharp}.
The finite-size comparisons described a sharper majority response on
larger grids, but the estimated midpoint continued to depend on the
observation period (Fig.~\ref{fig:finite-size}; Supplementary
Table~\ref{tab:revision-horizons}). A market remaining mixed at one
endpoint could therefore require further observation before its later
outcome was known. Neither the selected images nor the sampled crossing
curves established a threshold for indefinitely large or long-lived
markets~\cite{binder1981finite,landau2021guide}. The resemblance to
droplets forming and joining also did not establish classical nucleation.
We measured no energy barrier, surface tension, formation rate, or
critical group size separating growth from decay. The supported
mechanism was the propagation of changed responses through neighboring
decisions over finite grids and observation windows.

The corrected execution mechanism changed the financial interpretation
of limited dealer resources. Dealer models commonly connect trading
to intermediary cash and inventory~\cite{garman1976market,stoll1978supply,ho1981optimal,braunmunzinger2018agent},
but holdings alone do not specify how those resources translate into
prices. Here, the constant-product reserve curve determined both the
uniform price that supported a batch and the marginal quote left after
settlement~\cite{angeris2022constant}. When the dealer bought shares,
its share reserve increased and its cash reserve decreased, so subsequent
net sales could execute at lower prices. Shallow positive reserves thus
allowed large proportional declines without requiring a large fraction
of eligible orders to be rejected (Fig.~\ref{fig:reservoir}C,D).
Zero depth was a separate fixed-price control, because the dealer took
no net position and only opposing trader orders matched. This distinction
prevented the zero-depth case from being read as evidence that every
small dealer protected the market from a crash. The no-switch comparisons
also mattered because shallow markets had large drawdowns while
retaining the original decision rule
(Supplementary Table~\ref{tab:paired-switch}). The paired increases
separated that baseline financial variation from the intervention.
The evidence therefore linked crash magnitude to the execution and
price mechanism as well as to available resources; resource scarcity
alone did not imply a small proportional decline.

Financial conditions could also feed back into the spread of selling.
Changing dealer depth or trader holdings altered completed orders and
prices, which then changed technical-trader decisions and the local
inputs seen by their neighbors (Figs.~\ref{fig:reservoir}
and~\ref{fig:balance-sheet}). Holding those traders prevented the same
financial changes from entering the interior update, providing an exact
control on this route of influence. With technical traders active,
majority probability did not consistently increase or decrease with
depth, even though drawdown at the comparison with $f=0.20$
decreased across positive depths. Freezing the Hold-band width tested
one contributor to feedback while preserving responses to a moving
reference price (Supplementary Table~\ref{tab:revision-band}). That
comparison did not remove technical trading, and its frozen widths
retained differences accumulated during each starting market's baseline.
The uncertainty in the paired contrast also prevented a clear attribution
of the depth difference to width adaptation.
The additional reserve-depth series changed funding and initial price
sensitivity together, so it could not identify their independent effects
(Supplementary Table~\ref{tab:revision-impact}). This coupling distinguishes
the experiment from models that assign a separate impact coefficient
to completed order imbalance~\cite{kyle1985continuous,farmer2005predictive,bouchaud2009how}.
Majority probability and drawdown consequently measured different
outcomes: the share of markets in which enough traders decided to sell,
and the size of the largest decline along the resulting price path.
Neither outcome alone specified the other.

The model differed from established agent-based markets in where traders
interacted and in what was changed. In the minority game, every trader reads
the same record of past outcomes, keeps a score for each of its strategies,
and uses the one that would have done best, so traders interact only through
the market outcome and the main control parameter is the number of
distinguishable histories relative to the number of
traders~\cite{challet1997minority,challet2000modeling}. In the Santa Fe
artificial stock market, traders learn forecasting rules from a shared price
and dividend~\cite{lebaron1999time}; adding technical rules to a trader's
repertoire was a dominant strategy that lowered wealth and raised volatility
for all~\cite{joshi1999technical}, and the rate at which traders revised
their rules settled at an equilibrium that left earnings below what slower
revision would have produced~\cite{joshi2002financial}. Models with
chartists and fundamentalists let traders move between the two groups at
rates set by recent profits or by the opinions of
others~\cite{lux1999scaling,brock1998heterogeneous,kirman1993ants}. Lattice
and cellular-automaton markets encoded local interactions through stochastic
imitation or spin-like responses, sometimes alongside market-wide
signals~\cite{bornholdt2001ising,iori2002microsimulation,bartolozzi2004stochastic,qiu2007cellular,sornette2006importance,fan2009effect,atman2012influence}.
Wolfram's one-dimensional market and a later model built on Rule~110 used
deterministic rules and set the price from the numbers of buyers and
sellers~\cite{wolfram2002new,valsan2022on}. These models explained aggregate market behavior through different forms
of learning and interaction. The present comparison focused instead on a
persistent change assigned to a specified subset of traders. Here
every trader kept a fixed rule, the intervention changed one entry of the
rule table for a chosen fraction of traders, and the outcome was compared
with a market that shared the same history. Only the technical traders on
the edge decided from the price, so the same rule change could be tested
with those traders active and held. Orders needed cash or shares and were
completed by other traders or by a dealer with limited holdings, so a Sell
majority and a price crash were measured separately. These choices gave up
the learning and strategic behavior of established models in exchange for a
rule space that can be enumerated and an intervention that can be repeated
exactly. The questions answered were correspondingly narrower: not whether
technical trading emerges or whether the market becomes efficient, but how
many traders had to change a rule before a Sell majority formed, and whether
its formation required feedback from price.

Several modeling choices limit how far these conclusions extend. Traders
followed assigned rules rather than choosing strategies to maximize returns,
and updated together on a regular grid under one baseline lookup table.
Different neighbor configurations could have the same sum and therefore
generate the same response. The technical traders followed one band rule
and occupied selected perimeter sites, leaving other placements and
price-based strategies untested. Unit orders cleared at a common price
against a constant-product dealer, a specific reserve mechanism rather
than a calibrated representation of an equity exchange. The model had
no order book, transaction fees, dividends, borrowing, or strategic dealer
risk management, and its time steps and price response were not fitted
to market data. Its simulated crashes therefore did not predict the
timing or magnitude of actual market crashes. Conservation and the
isolated-round-trip identity constrained settlement, but did not establish
economic realism or exclude profits earned from other traders. The
long no-switch controls extended observation of existing starting markets
without establishing their behavior indefinitely
(Appendix~\ref{app:revision-controls}). Likewise, neither the baseline
period nor an average over final states demonstrated that the distribution
of market states had stopped changing. Other decision rules, asynchronous
updates, alternative interaction networks, and different execution
mechanisms would test how the connection between local decisions and
price declines depended on these assumptions.
